\documentclass[10pt]{article}
\usepackage[T1]{fontenc}
\usepackage{lmodern,amsmath,amssymb,amsthm,microtype}
\usepackage[margin=0.82in]{geometry}
\usepackage[colorlinks=true,urlcolor=blue,citecolor=blue]{hyperref}
\newtheorem{theorem}{Theorem}
\newtheorem{corollary}{Corollary}
\newcommand{\tm}{t}
\title{A Binary Characterization of OEIS A195608}
\author{Qichao Wang\\{\small Hebei University of Technology}\\[0.3em]
Daoyu Dong\\{\small University of Electronic Science and Technology of China}}
\date{October 26, 2026}
\hypersetup{pdftitle={A Binary Characterization of OEIS A195608},pdfauthor={Qichao Wang and Daoyu Dong}}
\begin{document}
\maketitle
\begin{abstract}
We characterize every index at which the sum of the first $n$ evil
numbers equals one quarter of the product of the $n$th and $(n+1)$st
evil numbers. The binary digit parity of two adjacent integers decides
membership. This proves the three structural conjectures in OEIS
A195608: the gaps are exactly $1,4,6$, the maximal consecutive blocks
have length one or three, and each block start other than the initial
index has the stated neighboring-evil-number difference equal to five.
The initial index is an explicit boundary exception to the wording of
the third conjecture, since the referenced evil-number sequence starts
at index one.
\end{abstract}

\section{An exact membership criterion}
Let $\tm(k)$ be the parity of the number of ones in the binary
expansion of $k\ge0$. Then
\[
 \tm(2k)=\tm(k),\qquad \tm(2k+1)=1-\tm(k).
\]
Write $E(n)$ for the $n$th nonnegative evil number, using the
one-based indexing of OEIS A001969 \cite{evil}. Exactly one of $2k$
and $2k+1$ has even binary digit sum, so
\begin{equation}\label{eq:evil}
 E(n)=2(n-1)+\tm(n-1)\qquad(n\ge1).
\end{equation}
Let $\mathcal S$ consist of the positive integers $n$ satisfying
\begin{equation}\label{eq:def}
 \sum_{j=1}^n E(j)=\frac{E(n)E(n+1)}4.
\end{equation}
Its increasing enumeration is OEIS A195608 \cite{target}.

\begin{theorem}\label{thm:membership}
For every $k\ge0$,
\[
 2k+1\in\mathcal S\quad\Longleftrightarrow\quad \tm(k)=0.
\]
For every $k\ge1$,
\[
 2k\in\mathcal S\quad\Longleftrightarrow\quad
 \tm(k-1)=\tm(k)=0.
\]
Thus the entire set, with no search for the original summatory
identity, is
\[
 \mathcal S=\{2k+1:k\ge0,\ \tm(k)=0\}
 \;\cup\;\{2k:k\ge1,\ \tm(k-1)=\tm(k)=0\}.
\]
\end{theorem}
\begin{proof}
Put $T(n)=\sum_{j=0}^{n-1}\tm(j)$. Pairing consecutive terms gives
$T(2k)=k$ and $T(2k+1)=k+\tm(k)$. Equation~\eqref{eq:evil} gives
\[
 \sum_{j=1}^n E(j)=n(n-1)+T(n).
\]
For $n=2k+1$, put $u=\tm(k)$. Then
$\tm(n-1)=u$, $\tm(n)=1-u$, and subtracting four times the left
side of \eqref{eq:def} from its right-hand product gives
\[
 E(n)E(n+1)-4\sum_{j=1}^n E(j)=-2u.
\]
This vanishes precisely when $u=0$.

For $n=2k$, put $u=\tm(2k-1)=1-\tm(k-1)$ and
$v=\tm(2k)=\tm(k)$. The corresponding difference is
\[
 (4k-2)v+4ku+uv-4k.
\]
For $(u,v)=(0,0),(0,1),(1,0),(1,1)$ its values are,
respectively, $-4k,-2,0,4k-1$. Since $k\ge1$, only $(u,v)=(1,0)$
works, giving the stated condition.
\end{proof}

\section{All three structural questions}
OEIS A195608 lists three conjectures \cite{target}: its first
differences take only the values $1,4,6$; its maximal consecutive
blocks have size one or three; and at a block's first index $m$,
$E(m+1)-E(m-1)=5$. The last expression is defined only for $m\ge2$.
Our result resolves these statements with that necessary boundary
qualification.

\begin{corollary}\label{cor:blocks}
The maximal consecutive blocks in $\mathcal S$ have length one or
three. The differences of successive terms of its increasing
enumeration are exactly $1,4,6$. The initial block is $\{1\}$.
Every later block starts at an odd integer $m$ and satisfies
\[
 E(m+1)-E(m-1)=5.
\]
\end{corollary}
\begin{proof}
Every pair $\tm(2r),\tm(2r+1)$ consists of opposite bits. Therefore
neither $000$ nor $111$ occurs in the parity sequence. Consecutive
zero positions are at distance one, two, or three. A run of zero
positions has length at most two.

By Theorem~\ref{thm:membership}, a single zero at $k$ produces the
odd member $2k+1$. Two adjacent zeros at $k,k+1$ produce exactly
the three consecutive members $2k+1,2k+2,2k+3$. All even members
arise in this way. Thus all maximal blocks have length one or three.
Zero-position gaps of two and three produce member gaps four and
six; a zero-position gap of one is filled by its even member and
produces two gaps of one. All three differences occur, as the
initial terms $1,7,11,12,13$ demonstrate.

A noninitial block begins at $m=2k+1$ where $k\ge1$,
$\tm(k)=0$, and $\tm(k-1)=1$. Now
\[
 E(m+1)=4k+2+\tm(2k+1)=4k+3,
 \qquad
 E(m-1)=4k-2+\tm(2k-1)=4k-2.
\]
Their difference is five. At the initial block $m=1$, $E(0)$ is
outside the domain of A001969; no assertion involving it is needed.
\end{proof}

\section{Reproducibility and scope}
The accompanying \texttt{verify.py} independently generates evil
numbers by testing binary digit parity and tests the original
integer equality \eqref{eq:def}. It compares this computation with
the theorem for every index through one million, checks the complete
first 10,000 published terms, and checks all complete blocks and
block-start differences in the same interval. A second independent
generation uses the binary substitution $0\mapsto01$, $1\mapsto10$.
These finite audits validate the implementation; the proofs above
cover every positive index.

\begin{thebibliography}{9}
\bibitem{target} V. Shevelev, OEIS A195608, The On-Line Encyclopedia
of Integer Sequences, \url{https://oeis.org/A195608}, accessed
October 26, 2026.
\bibitem{evil} The OEIS Foundation, OEIS A001969,
\url{https://oeis.org/A001969}, accessed October 26, 2026.
\end{thebibliography}
\end{document}
